\SourceLeaf{169}
\Needspace{9\baselineskip}
\pdfbookmark[1]{[IV.] Bilanz zwischen den Verpflichtungen und den Mitteln der Professorenwitwenkasse zu Göttingen.}{witwen-iv}
\begin{center}
[IV.]\par\medskip
\textit{Bilanz zwischen den Verpflichtungen und den Mitteln der Professorenwitwenkasse zu Göttingen.}
\end{center}
\noindent\hspace*{\parindent}Auf den Wunsch hin, den mir die Universitätskirchendeputation vor einigen Monaten mitgeteilt hat, habe ich die Bilanz der Professorenwitwenkasse neu berechnet und diese Arbeit jetzt abgeschlossen. Ich werde meinen Bericht so aufbauen, dass ich zuerst die nötigen allgemeinen Erläuterungen vorausschicke, dann die Bilanzrechnung selbst in knapper, leicht übersichtlicher Form aufstelle, anschließend die einzelnen Rechnungsposten näher erläutere und zuletzt die Vorschläge für die unmittelbar bevorstehende Periode anschließe.
\begin{center}\rule{.26\textwidth}{.3pt}\end{center}
\noindent\hspace*{\parindent}Die Aufstellung der Bilanz für eine Anstalt wie unsere Witwenkasse muss sich offensichtlich auf einen bestimmten \textit{Zeitpunkt} beziehen. Dass ich dafür diesmal den 1. Oktober 1851 gewählt habe, wird keiner ausführlichen Rechtfertigung bedürfen. Die erste Bilanzrechnung war auf den 1. Oktober 1845 bezogen. Aus naheliegenden Gründen muss der Abstand zwischen zwei aufeinanderfolgenden Prüfungen eine volle Anzahl von Jahren umfassen, damit das Ergebnis deutlicher hervortritt. Sollte sich infolge der neuen Prüfung schließlich eine Änderung der bisherigen Vereinbarungen als angemessen erweisen, wird man den Beschluss offensichtlich viel lieber auf den neuesten Zustand als auf den vor einem Jahr bestehenden stützen wollen. Die jetzigen Statuten schreiben zwar alle fünf Jahre eine neue Revision vor; der Zeitpunkt, zu dem mich die Aufforderung erreichte, ließ aber keine andere Wahl mehr zu. Durch diese Verlängerung des fünfjährigen Zeitraums auf einen sechsjährigen ist außerdem nicht nur nichts verloren, sondern vielmehr eine noch etwas deutlichere Ausprägung der Zustandsänderung gewonnen worden.

Das Wesen der gesamten Bilanzrechnung, sowohl der jetzigen als auch der von 1845, besteht darin, dass einerseits die Verpflichtungen der Anstalt und andererseits ihre Mittel auf den gleichwertigen Kapitalwert zurückgeführt werden, und zwar für alle Zukunft, nicht für das nächste Jahr oder einige Jahre. Nur so lässt sich einer Anstalt, die nur zum kleinsten Teil auf Beiträgen und größtenteils auf ihrem Vermögensbesitz beruht und deren Teilnehmerzahl in den letzten Jahrzehnten so stark gewachsen ist, die Tragfähigkeit für alle Zukunft sichern.

Obwohl diesmal ebenso wie 1845 alle Rechnungen doppelt durchgeführt sind, nämlich nach einem Zinssatz von \(3\tfrac{1}{2}\) und einem von 4 Prozent, habe ich es für hinreichend gehalten, hier nur die Ergebnisse nach Ersterem anzugeben. Ein Teil des Vermögens bringt tatsächlich nur \(3\tfrac{1}{2}\) Prozent; bei einem anderen, jetzt höher verzinsten Teil ist eine Zinssenkung in nicht allzu ferner Zeit nicht unwahrscheinlich. Jedenfalls gebietet die Vorsicht, bei solchen Rechnungen immer einen etwas niedrigeren Zinssatz zugrunde zu legen, als gegenwärtig üblich ist.

Als Sterblichkeitstafeln habe ich, soweit die Rechnungen davon abhängen, auch diesmal die von \textsc{Brune} benutzt, die zuverlässigsten, die es überhaupt gibt.

Sämtliche Rechnungselemente, die sich nur aus den Erfahrungen bei der Witwenkasse selbst ableiten lassen und deshalb zuverlässiger werden, wenn diese Erfahrungen einen längeren Zeitraum umfassen, habe ich für die jetzige Bilanzrechnung neu bestimmt, indem ich die früheren Erfahrungen mit den neu hinzugekommenen verbunden habe. Bei den einzelnen Positionen wird dies näher angegeben werden.

Schließlich bemerke ich noch, dass bei allen Geldangaben Goldwährung gemeint ist. Die Originalrechnungen sind zwar durchweg auf Bruchteile eines Talers genau geführt, diese Bruchteile sind im vorliegenden Auszug aber weggelassen. Daraus erklären sich einige scheinbare kleine
\SourceLeaf{170}
Abweichungen bei den angegebenen Summen, die hier und da eine oder ein paar Einheiten betragen können.

\Needspace{210pt}
\begin{center}
\textit{Bilanzrechnung der Witwenkasse für den 1. Oktober 1851.}\par\smallskip
\footnotesize\setlength{\tabcolsep}{1pt}
\begin{minipage}[t]{.48\linewidth}
\centering
Verpflichtungen.\par
\begin{tabular}{l r}
&Taler\\
\shortstack[l]{Kapitalwert des festen\\Teils der Pensionen}&\\
\shortstack[l]{1. für die jetzt vorhandenen\\Witwen . .}&28786\\
\shortstack[l]{2. für Witwen und Waisen\\der jetzigen Mitglieder\\. . . . . . . . . . . . . . . . . . .}&58463\\
\shortstack[l]{3. für Witwen und Waisen\\der künftig beitretenden\\Mitglieder . . . . . . . . . . . .}&40712\\
\shortstack[l]{4. Kapitalwert des variablen\\Teils der Pensionen nach\\jetziger Regelung . .}&25715\\
\shortstack[l]{5. Kapitalwert der sonstigen\\Ausgaben . . .}&12907\\\cline{1-2}
\multicolumn{1}{r}{Summe Taler}&166581
\end{tabular}
\end{minipage}\hfill\vrule\hfill
\begin{minipage}[t]{.48\linewidth}
\centering
Mittel.\par
\begin{tabular}{l r}
&Taler\\
Kapitalwert&\\
\shortstack[l]{6. des Ertrags der Apotheke\\. . . . . . . . .}&28571\\
\shortstack[l]{7. der Beiträge der jetzigen\\Mitglieder . . .}&8961\\
\shortstack[l]{8. der Beiträge aller künftig\\eintretenden Mitglieder\\. . . . . . . . . . . . . . . . . . .}&13596\\
\shortstack[l]{9. Geldvermögen der\\Witwenkasse . . . . . .}&122290\\\cline{1-2}
\multicolumn{1}{r}{Summe Taler}&173419
\end{tabular}\par\medskip
Das Ergebnis der Bilanzrechnung ist also ein \textit{Überschuss} von 6837 Talern.
\end{minipage}
\end{center}
